\section{Introducción}

Los intercambios descentralizados basados en creadores automáticos de mercado (AMM) y
los mercados centralizados (CEX) representan mecanismos complementarios de
formación de precios. Su interacción mediante arbitraje y diferencias de liquidez
permite estudiar el seguimiento de precios y la respuesta a shocks. El aporte esperado
no es una teoría general de estabilidad de mercados, sino un conjunto de mecanismos
verificables en un simulador: qué parte de la corrección procede del arbitraje, cómo
cambia con capacidad limitada y cómo la profundidad de liquidez transforma una venta
de tamaño fijo. El modelo reducido se subordina a esa evidencia computacional y no se
presenta como equivalente automático a un mercado real.

El aporte esperado para la investigación es un experimento reutilizable y auditable,
junto con una evaluación de dónde falla su descripción reducida. Para traders, las
comparaciones controladas pueden ilustrar cómo los supuestos de capacidad y liquidez
modifican la exposición de ejecución; son análisis de escenarios, no señales validadas
ni estimaciones de costos de negociación reales.

\section{Problema, pregunta e hipótesis}

\subsection{Problema científico}

No está suficientemente determinado, para el modelo mínimo propuesto, bajo qué
combinaciones de capacidad de arbitraje, profundidad de liquidez, fricciones de
ejecución y comportamiento de los agentes se logra seguimiento de precios, persiste
una discrepancia o se amplifica una perturbación. Una simulación que solo grafique
estos efectos no constituye por sí sola una explicación de las condiciones que los
producen.

\subsection{Pregunta principal}

¿Cómo afectan la capacidad de arbitraje y la respuesta de los proveedores de liquidez
al error máximo de seguimiento y a la recuperación tras shocks controlados en un mercado
CEX/DEX simulado, y con qué precisión clasifica un mapa discreto reducido la contracción
en escenarios de simulación reservados?

\subsection{Hipótesis de trabajo}

El seguimiento depende del balance entre realimentación correctiva (arbitraje) y
mecanismos potencialmente amplificadores (presión de negociación y liquidez
adaptativa). La variación de estos mecanismos produce regiones de comportamiento
distinguibles con métricas experimentales; el mapa reducido sirve como diagnóstico
local y debe fallar explícitamente cuando sus hipótesis dejan de ser razonables.

\begin{description}[style=nextline]
  \item[H1a.] Bajo supuestos definidos de capacidad, latencia y costes, el arbitraje
    reduce las discrepancias posteriores a una perturbación aislada.
  \item[H1b.] La retirada proporcional de liquidez aumenta el impacto de una operación
    posterior de tamaño fijo; ciertas reglas adaptativas de LP amplifican shocks en
    comparación con un contrafactual de liquidez estática.
  \item[H1c.] Un mapa estimado en una región de operación declarada puede distinguir
    contracción de no contracción en escenarios ABM reservados. La confirmación debe
    especificar de antemano la regla de decisión, una referencia independiente y una
    mejora mínima aceptable; el diagnóstico actual de 3/8 no demuestra esa mejora.
\end{description}

\section{Objetivos}

\textbf{General.} Desarrollar y evaluar un marco computacional reproducible que cuantifique cómo la
capacidad de arbitraje y las respuestas de liquidez afectan el seguimiento de precios
tras shocks controlados, y contraste los límites predictivos de un mapa reducido con
simulaciones basadas en agentes reservadas para evaluación.

\begin{enumerate}
  \item Formalizar el estado, los mecanismos de transacción del AMM, el arbitraje y las
    políticas de la taxonomía de cuatro roles aprobada por el autor mediante ODD \cite{grimm2020odd}.
  \item Construir un mapa discreto reducido con banda de no arbitraje y variables
    observables en el ABM; evaluar contracción local y condiciones de estabilidad de la
    linealización solo en las regiones donde aplique.
  \item Implementar y verificar un simulador reproducible con reservas, inventarios,
    costes y eventos ordenados de forma explícita.
  \item Diseñar experimentos factoriales o de sensibilidad para shocks, intensidad de
    arbitraje, profundidad y respuesta de LP.
  \item Contrastar resultados teóricos y simulados, incluyendo regiones de desacuerdo,
    robustez estadística y límites de extrapolación.
\end{enumerate}
